\section{Temporal-Basis Notation and Diagnostic Conventions}
\label{app:temporal_basis_notation}

\begin{table}[h]
\centering
\small
\caption{Indices and objects in the temporal-basis apportionment model.}
\label{tab:temporal_basis_notation}
\begin{tabular}{lll}
\toprule
Symbol & Shape or range & Meaning \\
\midrule
\(k\) & \(1,\ldots,K\) & source-group index \\
\(b\) & \(1,\ldots,B\) & temporal-basis index \\
\(j=(k,b)\) & \(1,\ldots,J\) & source-major, basis-minor coefficient index \\
\(\Phi\) & \(T\times B\) & nonnegative source-activity basis \\
\(C\) & \(K\times B\) & nonnegative source--basis coefficients \\
\(\mathbf c\) & \(J=KB\) & source-major vectorization of \(C\) \\
\(H_\Phi^{\mathrm{lag}}\) & \(mT\times J\) & unprojected coefficient fingerprints \\
\(\widetilde H_\Phi\) & \(mT\times J\) & background-projected fingerprints \\
\(\widetilde{\mathbf h}_{kb}\) & \(mT\) & fingerprint for coefficient \(c_{kb}\) \\
\bottomrule
\end{tabular}
\end{table}

The constant-activity model uses \(B=1\), \(\phi_1(t)=1\), and
\(c_{k1}=\theta_k\).  In the general model, diagnostics are computed for the
\(J\) coefficient fingerprints.  Source-level activity trajectories are
reconstructed as
\[
\widehat\theta_k(t)=\sum_{b=1}^{B}\widehat c_{kb}\phi_b(t);
\]
fingerprint columns are not summed to create an artificial source-level
diagnostic vector.

For a source-level ambiguity summary, IASA reports the largest eligible
coherence between coefficients belonging to different sources and retains the
source--basis pair that attains it.  Coherence and ray distance are undefined
for pairs involving a coefficient with \(v_j\le\tau_v\); tables display these
entries as N/A, exclude them from pairwise extrema, and report the weak
coefficient separately.
